Autoregressive language models train on all tokens in parallel, yet must generate them one at a time.
Diffusion language models remove this limit by denoising many tokens in each forward pass \citep{nie2026large, ye2025dream7bdiffusionlarge, team2026diffusiongemma}.

How well this works depends on the sampler.
In the continuous domain, many improvements have been driven by improved sampling algorithms \citep{song2020denoising}.
The same is true for the language domain, where sampling order has been shown to have a dramatic effect on task performance \citep{kim2025train, ben2026accelerated, kim2026klass}.

Any diffusion language sampler makes two choices for each position in its denoising window.
The first is \emph{where} the next token comes from: the model's prediction, the token already at that position, or fresh noise.
The second is \emph{when} the position is committed to the KV cache, after which it can no longer change.
Existing samplers tie the two choices together.
Samplers for masked diffusion freeze a token the moment it is accepted \citep{nie2026large, ben2026accelerated}, so acceptance and commitment are the same event.
Block samplers sit at the other extreme: every token waits until the whole block has converged \citep{arriola2025blockdiffusioninterpolatingautoregressive, team2026diffusiongemma}.

We separate the two choices and study each extreme on DiffusionGemma \citep{team2026diffusiongemma}, a uniform diffusion model whose default sampler is the entropy-bound (EB) sampler of \citet{ben2026accelerated}.
Both extremes are costly.
\emph{Committing too early} removes self-correction: when accepted tokens are frozen, macro accuracy falls from 73\% to 20\% as the entropy bound is relaxed, while the unfrozen sampler stays near 70\%.
This self-correction is a property of the model rather than the sampler, as it persists when renoising is removed.
\emph{Committing too late} wastes compute: tokens converge in a largely left-to-right order, so a converged prefix sits idle in the window, and after ten denoising steps more than half of the slots are spent on tokens that are already final.

These two findings motivate the Rolling Entropy Sampler (RES).
RES uses the entropy bound for both choices: it accepts tokens as the EB sampler does, and it commits the accepted prefix of the window as soon as that prefix is at least $s$ tokens long.
Every position outside the committed prefix stays free to change.
The minimum shift $s$ sets how long a token stays in the window, and gives a second knob, beside the entropy bound, for moving along the speed-quality frontier.

In summary, our contributions are as follows:
\begin{itemize}
    \item We write samplers as a routing rule and a commitment rule (Section~\ref{sec:motivation}), which separates accepting a token from committing it.
    \item We show that DiffusionGemma self-corrects without renoising, and that freezing accepted tokens destroys its speed-quality frontier (Section~\ref{sec:too-early}).
    \item We show that the EB sampler converges left to right, leaving over half of the window idle late in each block (Section~\ref{sec:too-late}).
    \item We introduce RES (Section~\ref{sec:method}). Across seven benchmarks it raises tokens per forward pass by 32\% at an average cost of 1.9 points of accuracy, and it extends the frontier beyond what relaxing the entropy bound alone reaches (Section~\ref{sec:experiments}).
\end{itemize}






